% ================================================================
%  conteudo/revisao.tex
% ================================================================

\section{Revisão da Literatura}

\subsection{Microalgas como Matéria-Prima para Biocombustíveis}

As microalgas são microrganismos fotoautotróficos unicelulares capazes de
converter CO\textsubscript{2} e luz solar em biomassa com eficiência
fotossintética superior à da maioria das plantas terrestres. Em termos de
produtividade de lipídeos por área cultivada, estimativas indicam que as
microalgas podem produzir de 20 a 80 vezes mais óleo por hectare do que
as oleaginosas convencionais, como soja e palma (CHISTI, 2007). Essa
característica, aliada à capacidade de crescimento em águas residuárias e
na ausência de solo arável, posiciona as microalgas como matéria-prima
estratégica para a produção de biocombustíveis de segunda geração.

O acúmulo de lipídeos em microalgas é fortemente influenciado pelas
condições de cultivo, em especial pela disponibilidade de nitrogênio.
Sob privação de nitrogênio, diversas espécies redirecionam o fluxo de
carbono da síntese de proteínas para a síntese de triacilgliceróis (TAGs),
podendo atingir teores lipídicos de até 50--70\% da massa seca (HO et al.,
2014). No entanto, essa condição de estresse compromete a produtividade de
biomassa, de modo que a maximização da produção de lipídeos totais requer
estratégias de cultivo em duas fases: crescimento em condições nutricionais
plenas seguido de privação seletiva de nutrientes.

\subsection{\textit{Chlorella fusca}: Características e Potencial}

\textit{Chlorella fusca} é uma microalga verde de água doce pertencente ao
filo Chlorophyta, com morfologia esférica e diâmetro celular de 2 a 10~$\mu$m.
Embora menos estudada que \textit{Chlorella vulgaris}, sua congênere de maior
relevância industrial, \textit{C. fusca} demonstra crescimento robusto em
ampla faixa de pH (5,0--9,0) e tolerância a variações de temperatura e
irradiância, características favoráveis ao cultivo em ambientes não
controlados ou em efluentes industriais (SOUSA et al., 2020).

O perfil lipídico de \textit{C. fusca} é composto predominantemente por
ácidos graxos saturados e monoinsaturados na faixa de C16--C18 --
principalmente ácido palmítico (C16:0), ácido oleico (C18:1) e ácido
linoleico (C18:2) --, perfil compatível com os requerimentos estequiométricos
do processo HEFA para produção de hidrocarbonetos na faixa de combustível de
aviação (C8--C16) (PEARLSON et al., 2013).

\subsection{Efluentes Industriais como Meio de Cultivo}

A utilização de efluentes industriais como substrato para cultivo de
microalgas é uma estratégia que concilia a redução do custo do meio e o
tratamento biológico de resíduos. Efluentes da indústria de processamento
de alimentos, laticínios e de fermentação industrial contêm concentrações
apreciáveis de nitrogênio amoniacal, fósforo e micronutrientes, que podem
substituir parcialmente ou integralmente os sais de meio sintético (ABDEL-RAOUF
et al., 2012).

Estudos demonstram que \textit{Chlorella} sp.\ cultivada em efluente de
laticínio diluído pode atingir produtividades de biomassa comparáveis às
obtidas em meio sintético, com remoção simultânea de até 80\% do nitrogênio
e 90\% do fósforo do efluente (MATA et al., 2010). Esses resultados indicam
que a integração do cultivo microalgal a sistemas de tratamento de efluentes
é uma estratégia promissora tanto do ponto de vista ambiental quanto
econômico.

\subsection{Processo HEFA para Produção de Biojet Fuel}

O processo HEFA é uma rota tecnológica certificada para produção de SAF a
partir de óleos vegetais, gorduras animais e lipídeos microalgais. O processo
envolve duas etapas catalíticas principais realizadas em reatores de leito
fixo sob pressão elevada de hidrogênio (35--170 bar) e temperaturas entre
300 e 450~°C (PEARLSON et al., 2013).

A primeira etapa, a hidrodeoxigenação (HDO), remove os grupos carboxila e
carbonila dos triglicerídeos por meio de reações de hidrogenação,
descarboxilação e descarbonilação, convertendo-os em \textit{n}-parafinas
de cadeia longa (C15--C18). A estequiometria simplificada representativa
da HDO pode ser expressa pela Equação~\ref{eq:hdo}:

\begin{equation}
    \text{C}_{n}\text{H}_{2n-2}\text{O}_{2} + 3\,\text{H}_{2}
    \;\longrightarrow\;
    \text{C}_{n}\text{H}_{2n+2} + 2\,\text{H}_{2}\text{O}
    \label{eq:hdo}
\end{equation}

A segunda etapa, o hidrocraqueamento com isomerização catalítica, quebra
as cadeias longas em hidrocarbonetos de menor massa molar e as isomeriza,
produzindo os cortes de nafta (C5--C8), SAF (C8--C16) e diesel verde
(C16+). A distribuição entre os cortes é controlada pelo catalisador
(geralmente zeólita bifuncional com metal nobre), temperatura e pressão
de hidrogênio (ZHU et al., 2011). Para fins de cálculo de rendimento, a
conversão de H\textsubscript{2} pode ser estimada em aproximadamente
0,04 kg de H\textsubscript{2} por kg de óleo processado (PEARLSON et al.,
2013).

\subsection{Análise Orçamentária de Bioprocessos}

A análise de custos de produção de biocombustíveis microalgais é uma
etapa crítica para avaliar a viabilidade comercial da rota tecnológica.
Estudos de techno-economic analysis (TEA) indicam que os principais
componentes de custo são o cultivo (incluindo meio, energia e mão de obra),
a colheita e concentração da biomassa e a extração de lipídeos, que em
conjunto podem representar mais de 80\% do custo total de produção (MATA
et al., 2010). A substituição do meio sintético por efluente industrial
pode reduzir os custos com nutrientes em 20 a 50\%, dependendo da composição
do efluente e da diluição necessária (ABDEL-RAOUF et al., 2012).

